\section{Zadanie 2}
Niech $X, Y, Z$ będą niezależnymi zmiennymi o rozkładzie Poissona. Niech $E(X) = 1, E(Y) = 2, E(Z) = 3$. Oblicz $P(X+Y+Z \leq 2)$.

\subsection*{Rozwiązanie}
Skoro $X,Y,Z$ to niezależne zmienne o rozkładzie Poissona, to ich suma również jest zmienną Poissona z wartością oczekiwaną 6(na podstawie wykładu, jest to również w książce jako Lemat 5.2). A zatem:


\[
P(X+Y+Z \leq 2) = P(X+Y+Z = 0) + P(X+Y+Z = 1) + P(X+Y+Z = 2)
\]
Wynikiem będzie więc:
\[
\frac{e^{-6} \cdot 6^0}{0!} + \frac{e^{-6} \cdot 6^1}{1!} + \frac{e^{-6} \cdot 6^2}{2!} = 25e^{-6}
\]

\subsection*{Odpowiedź}
\[
\boxed{25e^{-6}}
\]
